\section{Conclusion}
\label{sec:conclusion}

We presented the first mechanistic analysis of integer interval membership and modular arithmetic representations in a pretrained transformer. Using linear probing, regression probes, and activation patching on \gptsmall, we found that the model stores rich numerical information---comparison (99.1\%), interval membership (83.8\%), and modular arithmetic (58\%)---in its residual stream, far above chance and shuffled-label controls. This information exists despite the model performing at or below chance on all three tasks behaviorally.

Three findings stand out. First, interval membership is genuinely harder to decode than single comparison, with a 15 percentage point gap consistent with the hypothesis that intervals require composing two boundary checks. Second, ring counting answers are decodable at 15$\times$ above chance despite 0\% behavioral accuracy---the most extreme behavior--representation dissociation reported to date. Third, activation patching identifies layers~8--9 as causally important for both comparison and interval tasks, aligning with prior circuit-level findings.

These results have practical implications. Behavioral evaluation alone dramatically underestimates what models internally compute. Probing should be a standard companion to behavioral metrics in evaluating model capabilities. The ``last-mile'' failure---good internal representations that do not reach the output---suggests that fine-tuning or improved prompting could unlock latent capabilities without changing core representations.

\para{Future work.} We plan to extend this analysis to larger models (GPT-2 Medium/Large, Llama~3) to test whether behavioral accuracy improves while internal representations remain similar. Circular and sinusoidal probes~\citep{levy2024language, kadlcik2025pretrained} may better capture the periodic structure of modular arithmetic. Path patching at the attention head and MLP level would provide finer-grained circuit identification. Finally, studying whether fine-tuning on these tasks closes the behavior--representation gap would directly test the ``last-mile'' hypothesis.
